\documentclass[sigplan,10pt]{acmart}

\setcopyright{none}
\renewcommand\footnotetextcopyrightpermission[1]{}
\settopmatter{printfolios=true}
\pagestyle{plain}

\acmConference[ATC '26]{2026 ACM SIGOPS Annual Technical Conference}{November 15--18, 2026}{Hyatt Hotel, Shatin, Hong Kong}
\acmYear{2026}
\acmISBN{}
\acmDOI{}

\title[ActPlane]{ActPlane: OS-Enforced Agent Harnesses}

\author{Paper \#XXX}
\affiliation{%
  \institution{Anonymous Submission}
  \city{}
  \country{}}

\begin{document}
\maketitle

Modern coding agents run inside AI agent harnesses: software around the model that
maintains the agent loop, routes tool calls, mediates resources, and returns
feedback. Their behavior is governed largely by natural-language instruction
files: \texttt{CLAUDE.md}, \texttt{AGENTS.md}, and similar files tell them to run
tests before committing, avoid direct pushes, keep secrets local, use review
tools, or ask for confirmation before destructive operations. These instructions
are useful but probabilistic. A cooperative agent can forget them during a long
task, or satisfy the wording at the tool layer while still reaching the same
effect through \texttt{bash -c}, a Python subprocess, a helper binary, or an
SDK. The result is a mismatch between the layer where policies are written and
the layer where agent actions actually happen.

Existing defenses do not close this gap. Prompt-only constraints have no
deterministic enforcement. Tool-call guardrails over tool
names and arguments are deterministic, but they only see effects that cross the
guarded tool API. Whole-agent sandboxes draw a coarse perimeter, but cannot
easily express workflow invariants such as ``commit only after tests'' or
``this derived data may leave only after redaction.'' The OS boundary is the
common layer below these routes: every exec, file access, and socket connect
crosses a kernel interface regardless of how the agent initiated it.

ActPlane is a programmable OS-level runtime policy enforcement system for AI
agent harnesses. It compiles declarative source/sink/gate policies into an
eBPF/BPF-LSM backend that performs
\emph{labeled information-flow control}: it maintains named labels over
processes, files, and endpoints. Labels propagate through
fork/exec, file reads and writes, and connect-level endpoint edges. Rules then check whether a
labeled subject may perform an action on a matching target. The policy language
is intentionally agent-oriented: it expresses secret-derived egress, untrusted
input requiring review, read-only sub-agents, mandatory mediation through a
tool, and temporal gates such as testing before commit.

The central design choice is to treat agent policies as information-flow
constraints rather than as tool-call filters. A source introduces a label, for
example a process lineage started by \texttt{actplane run}, a repository secret,
or a directory containing untrusted issue text. A sink names the operation to
guard, for example an executable path, a file path, or a network endpoint. A
gate records an event that justifies an otherwise forbidden flow, such as a
redaction tool, a review step, a successful test run, or an explicit
confirmation. A rule combines those objects with an explicit effect and a
human-readable reason. This lets the policy say what must be true about the
agent's execution history, not merely which high-level tool name appeared in a
transcript.

The key semantic distinction is that ActPlane does not pretend every backend
can implement every effect. \texttt{block} means pre-operation denial through a
BPF-LSM hook returning \texttt{-EPERM}; it is not supported in tracepoint mode.
\texttt{notify} reports while allowing the action to proceed. \texttt{kill}
terminates the violating task and is useful for harness-level failure when
pre-operation arguments are unavailable. Each rule declares its effect
explicitly; there is no global fallback that silently converts \texttt{block}
into \texttt{kill}.

This distinction matters for the paper's threat and harness model. ActPlane is
not trying to prove that an arbitrary malicious process cannot escape every
possible policy. It targets a common but underspecified workload: a powerful
coding agent that is mostly cooperative, has long-running context, can call
shells and helpers, and needs deterministic failure feedback when it crosses a
project boundary. In that model, OS-level enforcement is useful even when the
effect is harness-level termination rather than pre-operation denial. Killing a
violating exec after \texttt{sched\_process\_exec}, for example, is not the same
security primitive as LSM denial, but it is still an effective failed action
from the agent's perspective and can be paired with a specific remediation.

ActPlane also closes the feedback loop. Kernel events contain only rule
identifiers and enforcement metadata; the Rust runner maps those identifiers
back to rule names, reasons, and remediation text, appends a model-facing
payload to \texttt{.actplane/last-violation.txt}, and exposes a post-tool hook
adapter that feeds new violations into the next agent turn. This keeps the
kernel as the sole authority for matches while giving a cooperative agent enough
semantic context to choose an allowed path.

The implementation follows this model with a deliberately small kernel
interface. The BPF side normalizes hook-specific observations into a common
event record containing subject identity, action kind, target descriptor, label
masks, and the rule that fired. File and exec hooks perform target extraction;
network hooks extract endpoint descriptors; process hooks maintain lineage. The
common matcher then evaluates required and forbidden labels against compiled
rule tables. User space owns YAML parsing, glob lowering, rule explanation,
feedback-file management, and command launching. The default user experience is
therefore a wrapper command such as \texttt{actplane run <agent-or-command>},
with an optional policy file when the current directory's policy should not be
used.

The evaluation separates mechanism from user experience. In RQ1, we run a
190-trace decision-compliance benchmark sampled from 38 directive-derived policies. Each
policy has allowed-effect compliant, lookalike compliant, visible-violation,
script-mediated, and opaque-fixture traces, and each trace is executed under
prompt-filter, tool-regex, ActPlane without structured feedback, and full
ActPlane. Full ActPlane reaches 75.8\% Decision Compliance Rate, compared
with 48.4\% for prompt-filter, 45.3\% for tool-regex, and 53.7\% for the opaque
feedback ablation. The gain comes mainly from reducing false negatives on
script and opaque runtime effects, while remaining false positives reflect
generated-policy precision limits rather than invalid traces.

The remaining experiments measure both cost and end-to-end usefulness. The
overhead experiment measures per-event costs and deterministic whole-agent task
time; at 32 active no-hit rules, ActPlane adds 1.9\% on agent trace replay and
6.5\% on a Linux build. The repository-grounded experiment runs a selected
20-task OctoBench subset and improves official reward from 0.788 without
enforcement and 0.818 with tool-regex to 0.888 with ActPlane. Finally, the
corpus experiment maps real \texttt{AGENTS.md}, \texttt{CLAUDE.md}, and similar
instructions into observable, expressible, robust, and temptable categories
before deriving the policy/task benchmark.

The mechanism is intentionally conservative. Kernel provenance systems such as
CamQuery~\cite{pasquier2018camquery} already showed that cross-channel tracking
and prevention are possible; agent-policy systems such as
AgentSpec~\cite{wang2025agentspec} and Progent~\cite{shi2025progent} show the
need for explicit runtime rules; and Fides~\cite{costa2025fides} and
CaMeL~\cite{debenedetti2025camel} show that information-flow abstractions are
useful for agent safety. The paper's bet is that coding agents make this old
capability newly useful at the OS boundary: the policy objects are project
instructions, the subjects are agent lineages and sub-agents, the failures must
be understandable to a model, and the useful definition of enforcement includes
both pre-operation denial and deterministic harness feedback.

\bibliographystyle{ACM-Reference-Format}
\bibliography{references}

\end{document}
